% Piano di Qualifica
% da compilare con il comando pdflatex Piano_di_Qualifica_x.x.x.tex

% Dichiarazioni di ambiente e inclusione di pacchetti
% da usare tramite il comando % Dichiarazioni di ambiente e inclusione di pacchetti
% da usare tramite il comando % Dichiarazioni di ambiente e inclusione di pacchetti
% da usare tramite il comando \input{../../util/hx-ambiente}

\documentclass[a4paper,titlepage]{article}
\usepackage[T1]{fontenc}
\usepackage[utf8]{inputenc}
\usepackage[english,italian]{babel}
\usepackage{microtype}
\usepackage{lmodern}
\usepackage{underscore}
\usepackage{graphicx}
\usepackage{eurosym}
\usepackage{float}
\usepackage{fancyhdr}
\usepackage[table]{xcolor}
\usepackage{longtable}
\usepackage{chngpage}
\usepackage{grffile}
\usepackage[titles]{tocloft}
\usepackage{hyperref}
\hypersetup{hidelinks}

\usepackage{../../util/hx-vers}
\usepackage{../../util/hx-macro}
\usepackage{../../util/hx-front}

% solo se si vuole una nuova pagina ad ogni \section:
\usepackage{titlesec}
\newcommand{\sectionbreak}{\clearpage}

% stile di pagina:
\pagestyle{fancy}

% solo se si vuole eliminare l'indentazione ad ogni paragrafo:
\setlength{\parindent}{0pt}

% intestazione:
\lhead{\Large{\proj}}
\rhead{\includegraphics[keepaspectratio=true,width=50px]{../../util/hivex_logo2.png}}
\renewcommand{\headrulewidth}{0.4pt}

% pie' di pagina:
\lfoot{\email}
\rfoot{\thepage}
\cfoot{}
\renewcommand{\footrulewidth}{0.4pt}

% spazio verticale tra le celle di una tabella:
\renewcommand{\arraystretch}{1.5}

% profondità di indicizzazione:
\setcounter{tocdepth}{4}
\setcounter{secnumdepth}{4}


\documentclass[a4paper,titlepage]{article}
\usepackage[T1]{fontenc}
\usepackage[utf8]{inputenc}
\usepackage[english,italian]{babel}
\usepackage{microtype}
\usepackage{lmodern}
\usepackage{underscore}
\usepackage{graphicx}
\usepackage{eurosym}
\usepackage{float}
\usepackage{fancyhdr}
\usepackage[table]{xcolor}
\usepackage{longtable}
\usepackage{chngpage}
\usepackage{grffile}
\usepackage[titles]{tocloft}
\usepackage{hyperref}
\hypersetup{hidelinks}

\usepackage{../../util/hx-vers}
\usepackage{../../util/hx-macro}
\usepackage{../../util/hx-front}

% solo se si vuole una nuova pagina ad ogni \section:
\usepackage{titlesec}
\newcommand{\sectionbreak}{\clearpage}

% stile di pagina:
\pagestyle{fancy}

% solo se si vuole eliminare l'indentazione ad ogni paragrafo:
\setlength{\parindent}{0pt}

% intestazione:
\lhead{\Large{\proj}}
\rhead{\includegraphics[keepaspectratio=true,width=50px]{../../util/hivex_logo2.png}}
\renewcommand{\headrulewidth}{0.4pt}

% pie' di pagina:
\lfoot{\email}
\rfoot{\thepage}
\cfoot{}
\renewcommand{\footrulewidth}{0.4pt}

% spazio verticale tra le celle di una tabella:
\renewcommand{\arraystretch}{1.5}

% profondità di indicizzazione:
\setcounter{tocdepth}{4}
\setcounter{secnumdepth}{4}


\documentclass[a4paper,titlepage]{article}
\usepackage[T1]{fontenc}
\usepackage[utf8]{inputenc}
\usepackage[english,italian]{babel}
\usepackage{microtype}
\usepackage{lmodern}
\usepackage{underscore}
\usepackage{graphicx}
\usepackage{eurosym}
\usepackage{float}
\usepackage{fancyhdr}
\usepackage[table]{xcolor}
\usepackage{longtable}
\usepackage{chngpage}
\usepackage{grffile}
\usepackage[titles]{tocloft}
\usepackage{hyperref}
\hypersetup{hidelinks}

\usepackage{../../util/hx-vers}
\usepackage{../../util/hx-macro}
\usepackage{../../util/hx-front}

% solo se si vuole una nuova pagina ad ogni \section:
\usepackage{titlesec}
\newcommand{\sectionbreak}{\clearpage}

% stile di pagina:
\pagestyle{fancy}

% solo se si vuole eliminare l'indentazione ad ogni paragrafo:
\setlength{\parindent}{0pt}

% intestazione:
\lhead{\Large{\proj}}
\rhead{\includegraphics[keepaspectratio=true,width=50px]{../../util/hivex_logo2.png}}
\renewcommand{\headrulewidth}{0.4pt}

% pie' di pagina:
\lfoot{\email}
\rfoot{\thepage}
\cfoot{}
\renewcommand{\footrulewidth}{0.4pt}

% spazio verticale tra le celle di una tabella:
\renewcommand{\arraystretch}{1.5}

% profondità di indicizzazione:
\setcounter{tocdepth}{4}
\setcounter{secnumdepth}{4}

\usepackage{hyperref}
\hypersetup{hidelinks}

\version{0.0.0}
\creaz{28 dicembre 2016}
\author{\LS, \AZ}
\supervisor{verificatori}
\uso{esterno}
\dest{Prof. Tullio Vardanega, Zucchetti S.P.A.}
\title{Piano di Qualifica}

\renewcommand{\arraystretch}{1.5}
\setcounter{tocdepth}{4}
\setcounter{secnumdepth}{4}

\begin{document}
\maketitle
% diario delle modifiche per il piano di progetto
% da includere con % diario delle modifiche per il piano di progetto
% da includere con % diario delle modifiche per il piano di progetto
% da includere con \include{diario}

\begin{diario}
	1.0.1 & {\GG} (Responsabile) & 03/02/2017 & Correzione degli errori individuati nelle prime tre sezioni del documento \\ \hline
	1.0.0 & {\PB} (Responsabile) & 09/01/2017 & Approvazione documento \\ \hline
	0.1.0 & {\MM} (Verificatore) & 07/01/2017 & Verifica documento \\ \hline
	0.0.9 & {\PB} (Responsabile) & 05/01/2017 & Stesura Organigramma \\ \hline
	0.0.8 & {\LB} (Responsabile) & 05/01/2017 & Stesura Consuntivo di Periodo \\ \hline
	0.0.7 & {\LB} (Responsabile) & 04/01/2017 & Stesura Preventivo di Periodo \\ \hline
	0.0.6 & {\LB} (Responsabile) & 03/01/2017 & Inserimento Gantt Diagrammi delle Attività \\ \hline
	0.0.5 & {\PB} (Responsabile) & 29/12/2016 & Stesura Pianificazione delle Attività \\ \hline
	0.0.4 & {\PB} (Responsabile) & 28/12/2016 & Stesura Introduzione \\ \hline
	0.0.3 & {\LB} (Responsabile) & 27/12/2016 & Stesura Modello di sviluppo \\ \hline
	0.0.2 & {\PB} (Responsabile) & 27/12/2016 & Stesura Analisi dei Rischi \\ \hline
	0.0.1 & {\LB} (Responsabile) & 26/12/2016 & Stesura scheletro \\ \hline
\end{diario}

\begin{diario}
	1.0.1 & {\GG} (Responsabile) & 03/02/2017 & Correzione degli errori individuati nelle prime tre sezioni del documento \\ \hline
	1.0.0 & {\PB} (Responsabile) & 09/01/2017 & Approvazione documento \\ \hline
	0.1.0 & {\MM} (Verificatore) & 07/01/2017 & Verifica documento \\ \hline
	0.0.9 & {\PB} (Responsabile) & 05/01/2017 & Stesura Organigramma \\ \hline
	0.0.8 & {\LB} (Responsabile) & 05/01/2017 & Stesura Consuntivo di Periodo \\ \hline
	0.0.7 & {\LB} (Responsabile) & 04/01/2017 & Stesura Preventivo di Periodo \\ \hline
	0.0.6 & {\LB} (Responsabile) & 03/01/2017 & Inserimento Gantt Diagrammi delle Attività \\ \hline
	0.0.5 & {\PB} (Responsabile) & 29/12/2016 & Stesura Pianificazione delle Attività \\ \hline
	0.0.4 & {\PB} (Responsabile) & 28/12/2016 & Stesura Introduzione \\ \hline
	0.0.3 & {\LB} (Responsabile) & 27/12/2016 & Stesura Modello di sviluppo \\ \hline
	0.0.2 & {\PB} (Responsabile) & 27/12/2016 & Stesura Analisi dei Rischi \\ \hline
	0.0.1 & {\LB} (Responsabile) & 26/12/2016 & Stesura scheletro \\ \hline
\end{diario}

\begin{diario}
	1.0.1 & {\GG} (Responsabile) & 03/02/2017 & Correzione degli errori individuati nelle prime tre sezioni del documento \\ \hline
	1.0.0 & {\PB} (Responsabile) & 09/01/2017 & Approvazione documento \\ \hline
	0.1.0 & {\MM} (Verificatore) & 07/01/2017 & Verifica documento \\ \hline
	0.0.9 & {\PB} (Responsabile) & 05/01/2017 & Stesura Organigramma \\ \hline
	0.0.8 & {\LB} (Responsabile) & 05/01/2017 & Stesura Consuntivo di Periodo \\ \hline
	0.0.7 & {\LB} (Responsabile) & 04/01/2017 & Stesura Preventivo di Periodo \\ \hline
	0.0.6 & {\LB} (Responsabile) & 03/01/2017 & Inserimento Gantt Diagrammi delle Attività \\ \hline
	0.0.5 & {\PB} (Responsabile) & 29/12/2016 & Stesura Pianificazione delle Attività \\ \hline
	0.0.4 & {\PB} (Responsabile) & 28/12/2016 & Stesura Introduzione \\ \hline
	0.0.3 & {\LB} (Responsabile) & 27/12/2016 & Stesura Modello di sviluppo \\ \hline
	0.0.2 & {\PB} (Responsabile) & 27/12/2016 & Stesura Analisi dei Rischi \\ \hline
	0.0.1 & {\LB} (Responsabile) & 26/12/2016 & Stesura scheletro \\ \hline
\end{diario}
\tableofcontents
\newpage

\section{Introduzione}
	\subsection{Scopo del documento}
	Questo documento ha lo scopo di fissare le strategie di verifica e validazione che il gruppo {\hx} ha deciso di adottare per perseguire gli obiettivi di qualità di processo e di prodotto relativi al progetto {\proj}. Il presente documento si propone di descrivere l'approccio del gruppo alle diverse fasi di verifica per poter ottenere il miglior risultato auspicabile in termini di qualità. Per perseguire tali obiettivi e risultati occorre verificare continuamente le attività svolte in modo da individuare e correggere velocemente le anomalie, minimizzando così l'utilizzo di risorse e allo stesso tempo mantenendo la correttezza del prodotto.

	\subsection{Scopo del prodotto}
	Lo scopo del prodotto è di implementare un editor \gloss{UML} che consenta, a partire dai diagrammi realizzati dall’utente, di generare il codice dell'applicazione da essi descritta, appartenente ad un determinato dominio applicativo: i giochi da tavolo. 
Il prodotto permetterà all’utente di realizzare l’applicativo desiderato utilizzando due tipologie di diagrammi ispirati ai diagrammi delle classi e delle attività UML, modificati appositamente per guidarlo nella progettazione dell'applicazione e agevolare la successiva generazione del codice, consentendo così una maggior qualità e strutturazione del codice stesso e di conseguenza un prodotto finito migliore.
L'editor dovrà essere fruibile dall'utente attraverso un browser desktop idoneo all'utilizzo delle tecnologie \gloss{HTML5}, \gloss{CSS3} e \gloss{JavaScript}, e il frutto del generatore di codice sarà scritto in \gloss{Java}.

	\subsection{Glossario}
	Per rendere più semplice e chiara la comprensione e di evitare ogni ambiguità relativa al linguaggio e ai termini utilizzati nei documenti, viene allegato il Glossario v_1_0_0 nel quale vengono raccolte le spiegazioni di termini tecnici o ambigui, di acronimi o abbreviazioni. Per rendere più facile la lettura, i termini presenti nel glossario saranno marcati con il pedice.
	
	\subsection{Riferimenti}
		\subsubsection{Normativi}
		\begin{itemize}
			\item \textit{Norme di Progetto v1.00}
			\item \textit{Capitolato d'Appalto C6}:
			\\ \href{http://www.math.unipd.it/~tullio/IS-1/2016/Progetto/C6p.pdf}{http://www.math.unipd.it/~tullio/IS-1/2016/Progetto/C6p.pdf};
			\item \textit{Standard ISO/IEC 12207:2008:}
			\\ \href{http://ieeexplore.ieee.org/document/4475826/}{http://ieeexplore.ieee.org/document/4475826/};
			\item \textit{Standard ISO/IEC 15504:}
			\\ \href{https://en.wikipedia.org/wiki/ISO/IEC_15504}{https://en.wikipedia.org/wiki/ISO/IEC_15504};
			\item \textit{Standard ISO/IEC 9126:}
			\\ \href{https://it.wikipedia.org/wiki/ISO/IEC_9126}{https://it.wikipedia.org/wiki/ISO/IEC_9126};
			\item \textit{PDCA(Plan-Do-Check-Act):}
			\\ \href{https://it.wikipedia.org/wiki/Ciclo_di_Deming}{https://it.wikipedia.org/wiki/Ciclo_di_Deming}.
		\end{itemize}
		\subsubsection{Informativi}
		\begin{itemize}
			\item \textit{Analisi dei Requisiti v1_0_0;}
			\item \textit{Piano di Progetto v1.00;}
			\item \textit{Glossario v1.00;}
			\item \textit{Indice di Gulpease:}
			\\ \href{http://it.wikipedia.org/wiki/Indice_Gulpease}{http://it.wikipedia.org/wiki/Indice_Gulpease};
			\item \textit{Regole del Progetto Didattico:}
			\\ \href{http://www.math.unipd.it/~tullio/IS-1/2016/Dispense/L09.pdf}{http://www.math.unipd.it/~tullio/IS-1/2016/Dispense/L09.pdf};
			\item \textit{Slide Qualità del Software:}
			\\ \href{http://www.math.unipd.it/~tullio/IS-1/2016/Dispense/L10.pdf}{http://www.math.unipd.it/~tullio/IS-1/2016/Dispense/L10.pdf};
			\item \textit{Slide Qualità del Prodotto:}
			\\ \href{http://www.math.unipd.it/~tullio/IS-1/2016/Dispense/L11.pdf}{http://www.math.unipd.it/~tullio/IS-1/2016/Dispense/L11.pdf}.
		\end{itemize}
\newpage
\section{Visione generale della strategia di gestione della qualità}
	\subsection{Obiettivi di qualità}
	\subsection{Qualità di processo}
		\subsubsection{Standard adottato}
		\subsubsection{Procedure per il controllo della qualità di processo}
		\subsubsection{Obiettivi di qualità di processo}
			\paragraph{Pianificazione di Progetto e Processo di Verifica e Controllo }
				\subparagraph{Rispetto della pianificazione}
				\subparagraph{Rispetto del budget}
			\paragraph{Processo di costruzione del software}
				\subparagraph{Rispetto della produttività di codifica}
			\paragraph{Processo di test per la qualifica del software/sistema}
				\subparagraph{Implementazione corretta delle funzionalità previste}
			\paragraph{Processo di gestione della documentazione del software}
				\subparagraph{Leggibilità della documentazione}
			\paragraph{Processo di Verifica del Software}
				\subparagraph{Copertura codice}
	\subsection{Qualità di prodotto}
		\subsubsection{Standard adottato}
		\subsubsection{Procedure per il controllo delle qualità di prodotto}
		\subsubsection{Obiettivi di qualità del prodotto}
	\subsection{Scadenze temporali}
\newpage
\section{La strategia di gestione della qualità nel dettaglio}
	\subsection{Risorse}
		\subsubsection{Risorse necessarie}
			\paragraph{Risorse umane}
			\paragraph{Risorse hardware}
			\paragraph{Risorse software}
		\subsubsection{Risorse disponibili}
			\paragraph{Risorse umane}
			\paragraph{Risorse hardware}
			\paragraph{Risorse software}
	\subsection{Misure e metriche}
		\subsubsection{Misure}
		\subsubsection{Metriche per i processi}
			\paragraph{Schedule Variance}
			\paragraph{Budget Variance}
			\paragraph{Produttività di codifica}
			\paragraph{Test di Unità eseguiti}
			\paragraph{Test di Integrazione eseguiti}
			\paragraph{Test di Sistema eseguiti}
			\paragraph{Test di Validazione eseguiti}
			\paragraph{Indice Gulpease}
			\paragraph{Copertura Codice}
	\subsection{Metriche per il prodotto}
		\paragraph{Requisiti obbligatori soddisfatti}
		\paragraph{Requisiti desiderabili soddisfatti}
\newpage
\section{Test}

\end{document}